% SHARD-ID: LB-14-TWO-PLUS-ONE
% SHARD-TITLE: The triangle in two space dimensions: anyon labels and 2D memory
% SHARD-SOURCES: theory/anyon-label-index.md; theory/anyon-selection-hybrid.md;
%              theory/memory-index-2d.md; theory/boundary-2d.md;
%              docs/2p1-antecedents.md; theory/verdicts/w21-joint-critic.md
% SHARD-CLAIMS: A-INDEX-TC-fin, M-INDEX-2D-fin, A-INDEX-PEPS, SHAPE-FLAT,
%              M-INDEX-2D-spec

% --- shard-local semantic macros (see WRITING-GUIDE rule 10) ---------------
\newcommand{\tpLab}{\mathsf{A}}      % anyon label set; kept distinct from \A
\newcommand{\tpCat}{\mathcal{C}}     % the input unitary fusion category
\newcommand{\tpDeph}{\mathcal{D}}    % pinching / dephasing map of a PVM
\newcommand{\tpWin}{W}               % a finite 2D window (disk or annulus)

\section{The triangle in two space dimensions: anyon labels and 2D
memory}\label{sec:two-plus-one}

The one-dimensional story of the preceding sections has three corners: a
charge attached to an endpoint, a memory recorded by a windowed observable,
and a soft theorem relating the two. In two space dimensions the first two
corners change character in a way that is worth stating plainly before any
theorem. An endpoint is no longer labelled by a group element or a projective
class; it is labelled by an \emph{anyon type}, an object of the Drinfeld
centre of the input fusion category, and those labels do not in general
subtract. A memory is no longer the displacement of a point-like kink; it is a
ledger of charge crossing a closed curve, a disk or an annulus, whose boundary
grows with the window instead of consisting of two endpoints. This section
reports what the campaign proved about those two corners, and it is honest
about the third: no two-dimensional soft theorem is claimed here.

That honesty is not modesty for its own sake. A survey of the continuum
antecedents produced three warnings, and each of them is a live trap for a
lattice calculation in $2+1$ dimensions.

\emph{First, memory means something different in $2+1$.} Direct continuum
computations show that permanent-displacement memory is an accident of
$d=4$. In $d=3$ the retarded field carries a tail, the force decays only as
$U^{-1/2}$, and the verdict is an ``infinite momentum memory effect'' with no
gravitational memory at all. A lattice observable in $2+1$ dimensions should
therefore be a momentum or velocity kick, stated as a rate or on a finite
window, and never as a finite permanent displacement. If a lattice computation
returns a finite permanent shift, the honest first hypothesis is that the
lattice regulator produced it, not the physics. Everything below respects this:
the memory statements of this section are \emph{charge} ledgers on a finite
window, and no displacement is defined.

\emph{Second, a confined or dressed charge sector silently trivialises the
triangle.} Logarithmic confinement in three-dimensional quantum
electrodynamics forbids finite-energy asymptotic charged states, so the
asymptotic symmetry acts trivially on the scattering matrix and the memory
effect is abolished. The lattice analogue is immediate and dangerous: in a
gapped $2+1$-dimensional topological phase there is no radiative sector at
all, so a ``soft theorem'' derived there risks being the identity $0=0$ dressed
up in notation. The pre-registered falsifier is that the sector carrying the
charge must be shown to possess finite-energy asymptotic states --- gapless at
a boundary, or tuned to criticality --- before any status label is allowed to
move. The toric-code theorem below is stated as a \emph{label} theorem
precisely so that it is not mistaken for a radiative soft theorem: it is
exactly and nontrivially true in a gapped phase, and it makes no soft claim.

\emph{Third, name the order of the fall-off.} The continuum literature
separates a radiative $1/r^{d/2-1}$ sector from a Coulombic $1/r^{d-3}$
sector, and ties only the null piece to flux and hence to the soft sector.
That analysis is stated for $d\ge4$ and does not cover $d=3$, where the two
exponents invert: the Coulombic term $1/r^{0}$ does not decay at all and
dominates the radiative $1/r^{1/2}$. This structural oddity is flagged here
rather than quietly extrapolated, and no result below depends on a
radiative/Coulombic split.

What the antecedents \emph{do} leave open is the holonomy route. Colour memory
measures the change in a flat connection by parallel transport of test
charges: an observable that needs no propagating radiation and no permanent
displacement. Its lattice counterpart is the change in a string or matrix
product operator holonomy encircling a region, whose labels are exactly the
central idempotents of a tube algebra. That is the route taken by the first two
theorems below.

\subsection{The toric-code endpoint carries an exact anyon label}

\begin{definition}[The toric-code model, its boundary-circle measurement, and
the same-circle protocol]\label{def:toric-registers}
Let $\Lambda=C_{L_x}\,\square\,C_{L_y}$ with $L_x,L_y\ge3$ be a finite square
cellulation of the torus, with one qubit on each edge and Pauli operators
$X_j,Z_j$. For each vertex $v$ and plaquette $p$ put
\begin{equation}
  A_v:=\prod_{j\ni v}X_j,\qquad
  B_p:=\prod_{j\in\partial p}Z_j,\qquad
  H_{\mathrm{TC}}:=-\sum_v A_v-\sum_p B_p .
  \label{eq:tc-hamiltonian}
\end{equation}
A star and a plaquette share zero or two edges, so all these terms commute.
Let $\Omega$ be any torus ground vector, $A_v\Omega=B_p\Omega=\Omega$. The two
global relations $\prod_vA_v=\prod_pB_p=I$ leave a four-dimensional ground
space, but no choice of logical ground state enters anything below.

The \emph{anyon label set} is the pointed Drinfeld centre
\begin{equation}
  \tpLab:=\operatorname{Irr}Z(\operatorname{Vec}_{\Z_2})
   =\{1,e,m,\epsilon\}\;\cong\;\Z_2^{e}\times\Z_2^{m},
  \label{eq:tc-labels}
\end{equation}
written $x=(x_e,x_m)$ with $e=(1,0)$, $m=(0,1)$, $\epsilon=(1,1)$; addition in
\eqref{eq:tc-labels} is fusion. We write $\tpLab$ rather than $\mathcal A$ to
keep this label set distinct from the asymptotic symmetry group $\A$ of
\cref{sec:corner-a}. The group structure of \eqref{eq:tc-labels} is special to
this pointed model, and it is exactly what makes a difference of labels
meaningful.

Choose a proper contractible disk $D\subset\Lambda$ with compatible vertex and
plaquette sets $V(D)$, $P(D)$, and define the closed boundary operators and
their joint projection-valued measure
\begin{equation}
  F_e(D):=\prod_{v\in V(D)}A_v,
  \qquad
  F_m(D):=\prod_{p\in P(D)}B_p,
  \qquad
  P_a(D):=\tfrac14\bigl(I+(-1)^{a_e}F_e(D)\bigr)
                 \bigl(I+(-1)^{a_m}F_m(D)\bigr).
  \label{eq:tc-pvm}
\end{equation}
Every interior edge occurs twice in the appropriate product, so $F_e$ and
$F_m$ are closed lattice boundary operators; they commute, are self-adjoint,
and square to the identity. Equivalently, the two commuting bits
$\hat q_D=\bigl((I-F_e)/2,\,(I-F_m)/2\bigr)$ take values in
\eqref{eq:tc-labels}, $F_e$ reading the parity of enclosed electric charge and
$F_m$ that of enclosed magnetic charge.

Let $\gamma$ be a direct-lattice path and $\gamma^{*}$ a dual-lattice path,
each with one endpoint inside $D$ and one outside, so that each crosses
$\partial D$ an odd number of times. Put
\begin{equation}
  Z(\gamma)=\prod_{j\in\gamma}Z_j,
  \qquad
  X(\gamma^{*})=\prod_{j\pitchfork\gamma^{*}}X_j,
  \qquad
  W_x=Z(\gamma)^{x_e}\,X(\gamma^{*})^{x_m},
  \label{eq:tc-ribbon}
\end{equation}
with an irrelevant Pauli phase chosen so that $W_x$ is unitary. \emph{Pure}
means that the single $x$-component of \eqref{eq:tc-ribbon} has been selected;
a coherent sum of different $W_x$ need not have one label and lies outside the
definite-sector clause. For $x=\epsilon$ the direct and dual endpoints in one
small neighbourhood form the composite endpoint. For $x=1$, $W_1=I$ is the
empty ribbon and is not a physical endpoint.

The \emph{same-circle protocol} is: for an arbitrary initial density matrix
$\rho$, measure the projection-valued measure \eqref{eq:tc-pvm}, apply $W_x$,
and measure the \emph{same} measure again. Its joint law and its sector
increment are
\begin{equation}
  p_x(a,b):=\operatorname{Tr}\bigl[P_bW_xP_a\,\rho\,P_aW_x^{\dagger}P_b\bigr],
  \qquad
  \Delta q_D:=b-a\in\tpLab .
  \label{eq:tc-tpm}
\end{equation}
Using the \emph{same} finite boundary register twice is load-bearing, exactly
as in the finite windowed memory theorem of \cref{sec:memory-index}: the two
outcomes belong to one affine register before their difference is taken. Here
the register has zero offset and is the fusion group itself, and
\eqref{eq:tc-tpm} is not spectral arithmetic for a difference of noncommuting
operators.
\end{definition}

\provenance{A-INDEX-TC-fin}{\texttt{theory/anyon-label-index.md} \S\S1.1--1.4}%
{\texttt{theory/checks/anyon\_label\_check.py} ANYON-C0--C4 on a $4\times4$
torus (exact rational weights)}%
{\texttt{theory/verdicts/w21-joint-critic.md} \S(W2)}

\begin{theorem}[A toric-code ribbon endpoint has a definite anyon label and an
exact same-circle shift]\label{thm:toric-endpoint}
\statusProved\ In the setting of \cref{def:toric-registers}, for every
$a,b,x\in\tpLab$:
\begin{enumerate}
\item \emph{(Definite endpoint sector.)} The ribbon has exact boundary
  covariance,
  \begin{equation}
    F_eW_x=(-1)^{x_e}W_xF_e,
    \qquad
    F_mW_x=(-1)^{x_m}W_xF_m,
    \label{eq:tc-covariance}
  \end{equation}
  and hence the selection rule
  \begin{equation}
    P_bW_xP_a=\delta_{b,a+x}\,W_xP_a .
    \label{eq:tc-selection}
  \end{equation}
  In particular $P_bW_x\Omega=\delta_{b,x}W_x\Omega$: the enclosed string
  endpoint sits in the definite Drinfeld-centre sector $x$.
\item \emph{(Quantised two-measurement law.)} The law \eqref{eq:tc-tpm} is a
  probability distribution supported on $\{(a,b):b=a+x\}$, so every run has the
  group-valued outcome $\Delta q_D=x$. Starting from any torus ground state,
  $p_x(1,x)=1$.
\item \emph{(Braiding readout.)} A closed probe ribbon of type $y=(y_e,y_m)$
  around the circle may be represented as
  $L_y(D):=F_m(D)^{y_e}F_e(D)^{y_m}$ and obeys
  \begin{equation}
    L_yW_x=(-1)^{\,x_ey_m+x_my_e}\,W_xL_y ,
    \label{eq:tc-braiding}
  \end{equation}
  a nondegenerate pairing, so the two measured bits distinguish all four
  sectors.
\end{enumerate}
The theorem is finite-volume and exact; it assumes neither a thermodynamic
limit nor any tensor-network approximation.
\end{theorem}

\begin{scope}
The label shift is invariant under ribbon deformations that preserve the
inside/outside endpoint partition, including arbitrary extra even crossings of
the circle. It need \emph{not} survive moving an endpoint across $\partial D$,
and it need not survive changing the measured circle between the two
measurements: both fences are part of the statement. A coherent superposition
of different ribbon types is not a pure endpoint and is outside the theorem.
Nothing here is a soft theorem, a current identity, a Ward identity, or a
persistence statement; and nothing here alters or repackages any
one-dimensional row of the campaign. The finite windowed memory theorem of
\cref{sec:memory-index} is used only as the protocol design pattern --- the
same register measured twice --- and no current, wall coordinate, scattering
limit, or Gram inverse is imported from it.
\end{scope}

\emph{Idea of proof.} A direct-lattice path enters and leaves every internal
vertex, so it overlaps the corresponding star in zero or two edges, while at
either endpoint the overlap is a single edge; the Pauli relation $XZ=-ZX$ then
delivers exactly one minus sign per endpoint. The dual statement holds for
$X(\gamma^{*})$ against the plaquettes. Multiplying the star identities over
$V(D)$ and the plaquette identities over $P(D)$ cancels every interior
incidence modulo two and leaves exactly one endpoint sign in each active
component, which is \eqref{eq:tc-covariance}. Moving $F_s$ through $W_x$ in
each spectral factor of \eqref{eq:tc-pvm} gives $P_bW_x=W_xP_{b-x}$, and right
multiplication by $P_a$ together with orthogonality of the measure gives
\eqref{eq:tc-selection}. On a ground vector every factor of $F_e$ and $F_m$
has eigenvalue $+1$, so $P_a\Omega=\delta_{a,1}\Omega$ and the definite-sector
clause follows. For the second clause, each weight in \eqref{eq:tc-tpm} is the
trace of $K_{a,b}\rho K_{a,b}^{\dagger}$ with $K_{a,b}=P_bW_xP_a$, hence
nonnegative; it vanishes off the support because the Kraus operator itself
vanishes there, before any property of $\rho$ is used; and summing over $b$
first, using $\sum_bP_b=I$, cyclicity, and $W_x^{\dagger}W_x=I$, gives total
mass $\operatorname{Tr}\rho=1$ with no assumption that $\rho$ commutes with the
first measurement. For the third clause, the closed $Z$-type boundary string
$F_m$ has odd intersection with the active dual $X$-string exactly when
$x_m=y_e=1$, and the direct/dual exchange of that argument gives the other
factor; every nontrivial $x$ is detected by the probe $y=m$ when $x_e=1$ and
by $y=e$ otherwise.

There is also an exact non-abelian instance, which is worth recording because
it shows precisely why the group-valued reading is special. For a Levin--Wen
model with Ising input, whose bulk category is the doubled Ising theory, take
the endpoint type and the enclosed source both equal to $(\sigma,\bar1)$.
The exact bulk fusion rule is
$(\sigma,\bar1)\otimes(\sigma,\bar1)=(1,\bar1)\oplus(\psi,\bar1)$, so a
boundary-circle measurement may return either of two outcomes, each through a
one-dimensional channel. The retained datum is the fusion event
$(x,a\to b;\mu)$, not a quotient $b/a$; the two possible outputs are the
explicit obstruction to subtracting Drinfeld-centre names. A group-valued
shadow survives --- the universal grading of Ising is $\Z_2$, giving
$U(Z(\mathrm{Ising}))\cong\Z_2\times\Z_2$, in which both outcomes have the same
grade and $|b|-|a|=|x|$ --- but that shadow conserves only a coarse parity and
chooses neither channel. This calculation is exact fusion arithmetic; its
identification with a chosen microscopic endpoint tensor is \emph{not} proved,
and is not part of \cref{thm:toric-endpoint}.

\provenance{A-INDEX-TC-fin}{\texttt{theory/anyon-label-index.md} \S2 (theorem)
and \S3 (proof, steps (1.1)--(1.4)); the doubled-Ising instance is \S4 and is
not part of the promoted row}%
{\texttt{theory/checks/anyon\_label\_check.py} ANYON-C0--C4, with the
registered red mode \texttt{-{}-red wrong-sector}; fusion arithmetic separately
in \texttt{theory/checks/cat\_hunt\_check.py} CATH-C1--C2}%
{\texttt{theory/verdicts/w21-joint-critic.md} \S(W2) and its promotion table}

\subsection{Two-dimensional memory is a charge ledger on a finite window}

\begin{definition}[Two-dimensional windows, the background-subtracted window
charge, and its two-measurement law]\label{def:2d-window-charge}
Let $\Gamma=\Z^{2}$, or any countable locally finite two-dimensional lattice,
with uniformly finite on-site Hilbert spaces and quasi-local algebra $\alg$.
Let a compact group act site by site, choose one closed one-parameter subgroup
and a Hermitian on-site generator $q_x$, and impose the integral-circle
condition of \cref{sec:memory-index} in the form
\begin{equation}
  e^{2\pi i q_x}=c\,I,\qquad c=e^{2\pi i\kappa},\quad\kappa\in[0,1),
  \qquad\text{equivalently}\qquad
  \spec q_x\subset\kappa+\Z .
  \label{eq:2d-int}
\end{equation}
Compactness of the group alone is not silently turned into a charge theorem: a
circle subgroup and \eqref{eq:2d-int} must be selected, and a finite or
disconnected symmetry with no such generator is outside every statement below.
Let $\alpha_t$ be the infinite-volume dynamics; the finite theorem needs only
that it is a group of $C^{*}$-automorphisms.

The declared protocol windows are disks and annuli,
\begin{equation}
  B_R(z)=\{x:\dist(x,z)\le R\},
  \qquad
  A_{r,R}(z)=\{x:r<\dist(x,z)\le R\},
  \label{eq:2d-windows}
\end{equation}
both finite; a disk has one boundary component and an annulus two, each
oriented outward from the annular region. The finite theorem in fact applies to
every finite $\tpWin\subset\Gamma$. A disk exhaustion $B_{R_m}(z)$ exhausts the
plane; an annular exhaustion $A_{r_0,R_m}(z)$ with fixed $r_0$ exhausts the
exterior of a fixed core and not the plane, and nothing is called an
exhaustion here unless the domain it exhausts is named.

The measured operator is a finite sum with a scalar counterterm,
\begin{equation}
  \Qhat_{\tpWin}=\sum_{x\in \tpWin}q_x-b_{\tpWin}I .
  \label{eq:2d-window-charge}
\end{equation}
Two preparations are intended. For a \emph{symmetry-broken pair}, choose two
zero-temperature vacua, a reference interface curve, and the phase profile on
its two sides, and let $b_\tpWin$ be the corresponding sum of vacuum
densities; the curve fixes only the normal-ordering convention and no sharp
interface-position operator is assumed. For a \emph{symmetric vacuum}, take
$b_\tpWin=\abs{\tpWin}\rho_0$ with $\rho_0$ the vacuum density; there is then
no wall coordinate at all, and the protocol measures charge emitted from or
absorbed by the window. The scalar $b_\tpWin$ may be non-integral and may
change with $\tpWin$; this is harmless below precisely because both
measurements use the same window and hence the same offset.

Let $E_{\tpWin,t}$ be the spectral resolution of
$\Qhat_\tpWin(t)=\alpha_t(\Qhat_\tpWin)$. Measuring projectively at
$t_-<t_+$, the recorded escaped charge and its law are
\begin{equation}
  \nu:=q_--q_+,
  \qquad
  p_{\tpWin;t_-,t_+}(\nu)
   :=\sum_q\norm{E_{\tpWin,t_+}(\{q-\nu\})\,E_{\tpWin,t_-}(\{q\})\,\Psi}^{2},
  \label{eq:2d-tpm}
\end{equation}
absent spectral values contributing zero. This is a sequential measurement
law, not the spectral law of the generally noncommuting difference
$\Qhat_\tpWin(t_-)-\Qhat_\tpWin(t_+)$.
\end{definition}

\provenance{M-INDEX-2D-fin, D26}{\texttt{theory/memory-index-2d.md}
steps (1.1)--(1.4)}{\texttt{theory/checks/memory\_index\_2d\_check.py}
M2D-C1--C4}{\texttt{theory/verdicts/w21-joint-critic.md} \S(W3)}

\begin{theorem}[Every finite disk or annulus records an integer]%
\label{thm:memory-2d-finite}
\statusProved\ Assume \eqref{eq:2d-int}, a finite disk or annulus $\tpWin$, a
real scalar background $b_\tpWin$, a unit vector $\Psi$ in any representation
of $\alg$, and finite times $t_-<t_+$. Put
$\kappa_\tpWin:=\abs{\tpWin}\kappa-b_\tpWin \bmod \Z$. Then
\begin{equation}
  \spec\Qhat_\tpWin(t)\subset\kappa_\tpWin+\Z
  \qquad\text{for every automorphic time }t,
  \label{eq:2d-coset}
\end{equation}
and the law \eqref{eq:2d-tpm} is a probability distribution whose escaped
increment $\nu$ lies in $\Z$. No commutativity of the two Heisenberg
observables is assumed.
\end{theorem}

\begin{scope}
This is a finite, additive, $0$-form charge theorem and nothing more. It
implies no displacement: the one-dimensional conversion of escaped charge into
a unique kink position does not transplant, because a two-dimensional
interface has shape modes and a swept area rather than one canonical
coordinate. It implies no $1$-form or topological charge, no anyon content, no
loop-operator statement, and no result about topological order --- the
symmetry here is an ordinary on-site $0$-form symmetry. It supplies no channel
inventory and no ordered limit; the limit statement is
\cref{thm:memory-2d-limit} and is conditional. It uses no spectral gap and no
exponential clustering, so it holds verbatim in a gapless
symmetry-broken phase with Goldstone modes: gaplessness can change
probabilities, relaxation, and tightness, but it cannot move a same-window
branch off $\Z$ while \eqref{eq:2d-int} holds. Finally, the \emph{total}
unnormalised increment is what is integral; dividing by the perimeter
$\abs{\partial \tpWin}$ would give a different, density-valued law supported in
$\abs{\partial \tpWin}^{-1}\Z$, and that must never be advertised as the
integer-valued index.
\end{scope}

\emph{Idea of proof.} Each on-site eigenvalue differs from $\kappa$ by an
integer. The operators $\{q_x:x\in\tpWin\}$ commute because they act on
different sites, so finite spectral addition on \eqref{eq:2d-window-charge}
puts every spectral value of $\Qhat_\tpWin$ in
$\abs{\tpWin}\kappa-b_\tpWin+\Z$, and $\alpha_t$ preserves spectrum, giving
\eqref{eq:2d-coset} at every time with the same offset. Every summand of
\eqref{eq:2d-tpm} is a squared norm, and summing first over the final and then
over the initial resolution gives total mass one by two applications of
Parseval, with no interchange of the two noncommuting projections. Since the
observable, the window, and the background are the same at both times, both
$q_-$ and $q_+$ lie in the single coset $\kappa_\tpWin+\Z$, so
$\nu=q_--q_+\in\Z$ before any time or spatial limit. Every leaf of this
argument sees only one fixed finite set and its cardinality: none sees spatial
dimension, boundary length, a spectral gap, or a scattering channel, which is
why the one-dimensional proof transplants without loss.

What does change with dimension is the size of the boundary term, and it is
worth displaying because it is the obstacle to the limit theorem. For a
finite-range interaction $H=\sum_Z\Phi(Z)$ with a termwise charge-conserving
decomposition $[\Phi(Z),\sum_{x\in Z}q_x]=0$, every term lying wholly inside or
wholly outside the window cancels, so
\begin{equation}
  \frac{d}{dt}\Qhat_\tpWin(t)
   = i\,\alpha_t\Bigl(
       \sum_{Z\,:\,Z\cap \tpWin\ne\emptyset,\ Z\cap \tpWin^{c}\ne\emptyset}
       \bigl[\Phi(Z),\textstyle\sum_{x\in \tpWin}q_x\bigr]\Bigr),
  \qquad
  \Bigl\lVert\frac{d}{dt}\Qhat_\tpWin(t)\Bigr\rVert
   \le C_{\Phi,q,r}\,\abs{\partial_r \tpWin}.
  \label{eq:2d-perimeter}
\end{equation}
On a disk of radius $R$ the number of contributing terms is $O(R)$, and on an
annulus $O(r+R)$ because both boundary components contribute, in place of the
two-endpoint boundary of an interval. This perimeter growth is a real
structural difference, catalogued alongside the other two-dimensional boundary
findings in \cref{sec:negative-results}; a Lieb--Robinson estimate localises
the terms but does not by itself cancel the perimeter factor.

\provenance{M-INDEX-2D-fin, D26, M-INDEX-fin}%
{\texttt{theory/memory-index-2d.md} steps (1.5)--(1.6), with the leaf-by-leaf
transplant ledger at step (1.6) and the perimeter audit at step (1.11)}%
{\texttt{theory/checks/memory\_index\_2d\_check.py} M2D-C2 and the registered
red mode \texttt{-{}-red fractional-charge}; donor one-dimensional
memory-index and common-subsequence gates}%
{\texttt{theory/verdicts/w21-joint-critic.md} \S(W3): promotable at proved
strength}

\subsection{Categorical selection: what a general anyon endpoint obeys}

\begin{definition}[The tensor-network typing hypotheses for an annular
endpoint register]\label{def:peps-typing}
Let $\tpCat$ be a finite unitary fusion category and let a finite,
matrix-product-operator-injective tensor-network lattice model supply all of
the following \emph{exact} data. These are hypotheses, not consequences.

\smallskip
\noindent\textbf{(PT1) Tube measurement.} An annular representation of the
tube algebra $\operatorname{Tub}(\tpCat)$ whose minimal central idempotents
$\{P_a\}_{a\in\operatorname{Irr}Z(\tpCat)}$ are mutually orthogonal and sum to
the identity on the chosen boundary-circle register.

\noindent\textbf{(PT2) Pure endpoint.} A small-circle resolution of the
endpoint tensor $T_x$ with $P^{\mathrm{end}}_yT_x=\delta_{y,x}T_x$ for one
simple $x\in\operatorname{Irr}Z(\tpCat)$. A sum over $x$ is not called a pure
endpoint.

\noindent\textbf{(PT3) Pulling-through module action.} Exact zipper,
associator, and pulling-through equations identifying the large-circle action
of $T_x$ on a source sector $a$ with the semisimple decomposition
\begin{equation}
  x\otimes a\;\cong\;\bigoplus_b\Hom(b,x\otimes a)\otimes b .
  \label{eq:pt-fusion}
\end{equation}

\noindent\textbf{(PT4) Protocol instrument.} The physical charge-creation
operation is supplied either as a normalised family of channel-resolved Kraus
maps $T_{(x,\mu)}$, or else the experiment is explicitly conditioned on the
success of one displayed Kraus map. Category data alone supply neither a
normalisation nor probabilities, and (PT4) is the only normalisation input
anywhere in this subsection.

\smallskip
For the momentum form, let $T_x(y;w)$ be a translated, windowed representative
of the same pure endpoint data, with endpoint position $y$ along a chosen
direction and all microscopic smearing data collected in $w$. On a finite
periodic boundary of length $L$ set
\begin{equation}
  T_x(k;w)=\sum_{y=0}^{L-1}e^{iky}\,T_x(y;w),
  \qquad k\in(2\pi/L)\Z .
  \label{eq:pt-fourier}
\end{equation}
In an infinite-volume protocol, ``a limit'' means a declared sequence of
windows and momenta in a declared operator, weak, or matrix-element topology
in which multiplication by the fixed finite-sector projectors is continuous.
\end{definition}

\provenance{A-INDEX-PEPS}{\texttt{theory/anyon-selection-hybrid.md} \S\S1--2;
canonical typing in \texttt{theory/anyon-label-index.md} \S5}%
{\texttt{theory/checks/cat\_hunt\_check.py} CATH-C0--C4;
\texttt{theory/checks/soft\_2d\_hunt\_check.py} S2DH gates}%
{\texttt{theory/verdicts/w21-joint-critic.md} \S(W6) and \S9 item 5}

\begin{theorem}[Fusion selection for an annular endpoint
register]\label{thm:fusion-selection}
\statusProvedConditional\ Assume (PT1)--(PT3) of \cref{def:peps-typing}, and
for the momentum form the notation and declared topology stated there; assume
(PT4) only for the instrument-normalisation clause. Then on the finite annular
register
\begin{equation}
  P_b\,T_{(x,\mu)}\,P_a=0
  \qquad\text{unless}\qquad
  \mu\in\Hom(b,x\otimes a)\ \text{ and }\ N_{xa}^{b}>0 ,
  \label{eq:fusion-selection}
\end{equation}
and consequently the same-circle instrument weight
\begin{equation}
  p_x(a,b,\mu)
   =\operatorname{Tr}\bigl(P_bT_{(x,\mu)}P_a\,\rho\,
                           P_aT_{(x,\mu)}^{\dagger}P_b\bigr)
  \label{eq:fusion-weight}
\end{equation}
is supported only on the fusion events allowed by \eqref{eq:fusion-selection},
normalised or postselected exactly as stipulated by (PT4). Moreover the
windowed/momentum form obeys the same rule,
$P_bT_x(k;w)P_a=0$ unless $\Hom(b,x\otimes a)\ne0$, at every finite momentum
and window, and every existing limit point in the declared topology obeys the
same zero-block constraint.

The retained fine datum is the fusion event $(x,a\to b;\mu)$. If $x$ is
invertible, tensoring by $x$ is an equivalence and the target is uniquely
$b=x\otimes a$, so the supplied index is group-valued,
$x\in\Inv(Z(\tpCat))$; the toric-code theorem
\cref{thm:toric-endpoint} is the case in which every simple object is
invertible. For general $x$ the only group-valued statement is the universal
grading shadow: presenting the universal grading group by $[c]=[r][s]$
whenever $N_{rs}^{c}>0$, every allowed event obeys
\begin{equation}
  \abs{b}=\abs{x}\abs{a},
  \qquad\text{equivalently}\qquad
  \abs{b}\abs{a}^{-1}=\abs{x}.
  \label{eq:grading-shadow}
\end{equation}
\end{theorem}

\begin{scope}
Three fences are part of this statement. First, no microscopic instantiation
beyond the toric code is claimed: supplying (PT2)--(PT4) for a particular
lattice model is an open, model-specific obligation, and it is not asserted
that every representative of a topological phase supplies them. Second, no
Taylor coefficient inside an allowed channel is claimed. The theorem says a
forbidden block vanishes; it says nothing about the value, the leading
behaviour, or the softness of an allowed block, and it therefore supplies no
categorical soft theorem and no universal soft coefficient. Third, this is one
selection-rule statement and is not an omnibus packaging of the group-label,
windowed-memory, and tensor-network results, which consume different
hypotheses and remain separate rows.

Two further limitations are worth naming. The labels used are \emph{bulk}
Drinfeld-centre labels, because the boundary circle sits in the
two-dimensional bulk; a physical edge or condensate requires its own module or
boundary-excitation category, under which bulk anyons may condense, be
identified, or split. And tube projectors select sectors only: they supply no
current operator, no Ward identity, no Gram inverse, no transition amplitude,
and no persistence theorem. The shadow \eqref{eq:grading-shadow} does not
select $b$, does not fix a probability, and does not authorise subtraction in
$\operatorname{Irr}Z(\tpCat)$; the doubled-Ising branching recorded after
\cref{thm:toric-endpoint} is the explicit instance.
\end{scope}

\emph{Idea of proof.} By (PT2) the endpoint carries one simple type $x$, and by
(PT3) its action on the source sector $a$ is the semisimple decomposition
\eqref{eq:pt-fusion}, which has no summand labelled $b$ when
$\Hom(b,x\otimes a)=0$. By (PT1) the $P_b$ are mutually orthogonal central
idempotents resolving the annular register, so left multiplication kills every
component of $T_x P_a$ except its $b$ summand; resolving the multiplicity
space, whose dimension is $N_{xa}^{b}$, gives \eqref{eq:fusion-selection}. A
forbidden block makes the corresponding Kraus operator in
\eqref{eq:fusion-weight} zero, so its weight vanishes by positivity alone;
normalisation is used only in the form supplied by (PT4). For the momentum
form, each summand $P_bT_x(y;w)P_a$ of \eqref{eq:pt-fourier} is zero on a
forbidden pair, and Fourier summation and windowing are linear, so their finite
linear combination vanishes. Finally, a net of identically zero blocks has zero
limit whenever multiplication by the fixed projectors is continuous in the
declared topology --- which asserts neither existence nor uniqueness of any
soft limit.

\provenance{A-INDEX-PEPS}{\texttt{theory/anyon-selection-hybrid.md}
step (1.1), leaves (2.1)--(2.6)}%
{\texttt{theory/checks/anyon\_label\_check.py} ANYON-C0--C4;
\texttt{theory/checks/cat\_hunt\_check.py} CATH-C0--C4;
\texttt{theory/checks/soft\_2d\_hunt\_check.py} S2DH gates. These finite gates
establish neither (PT2)--(PT4) for an arbitrary model nor the existence of an
infinite-volume momentum limit}%
{\texttt{theory/verdicts/w21-joint-critic.md} \S(W6) and \S9 item 5
(promoted after the structured-proof repair, orchestrator-verified)}

\begin{theorem}[Contractible deformations do not change the endpoint
morphism]\label{thm:isotopy-flat}
\statusProvedConditional\ Assume an exact fixed-point string or
matrix-product-operator representation together with its pulling-through and
associator equations. Let $\gamma$ and $\gamma'$ have the same resolved
endpoints and be related by a finite sequence of contractible local
pulling-through moves inside a disk containing no other insertion. Then, after
identifying the endpoint fusion spaces by the associator maps supplied by
those moves, the code-projected endpoint morphisms for $\gamma$ and $\gamma'$
agree.
\end{theorem}

\begin{scope}
The contractibility hypothesis is load-bearing and so is the fixed-point
hypothesis. A move sequence winding around another insertion is not
contractible in the complement of that insertion, so the coherence argument
does not identify it with the trivial path; its composition may instead be the
corresponding braid or monodromy action. No off-fixed-point statement is made:
nothing is claimed about a dressed string in a generic phase, and nothing is
claimed about a gapless on-shell soft factor.
\end{scope}

\emph{Idea of proof.} A single local pulling-through move is, by hypothesis, an
equality of the two code-projected tensor networks across that move, with the
endpoint fusion spaces identified by the associator map it supplies.
Composing those equalities along the finite move sequence identifies the first
network with the last, since equality and the supplied identifications are
stable under finite composition. Two contractible move sequences with the same
ends give the same categorical identification by pentagon coherence, which
relates their two finite compositions of associators. The exception is exactly
the non-contractible case, which the coherence argument does not reach.

\provenance{SHAPE-FLAT}{\texttt{theory/anyon-selection-hybrid.md}
step (1.2), leaves (2.1)--(2.5)}%
{\texttt{theory/checks/soft\_2d\_hunt\_check.py} with the registered
deformation mutation \texttt{-{}-red wrong-path}; the toric-code patch supplies
the nonzero witness (two paths give the same state, with an endpoint gap of
four)}%
{\texttt{theory/verdicts/w21-joint-critic.md} \S(W6) and \S9 item 5
(promoted, orchestrator-verified)}

\subsection{The conditional two-dimensional limit law}

\begin{definition}[Two-dimensional local-relaxation clauses and the charge
memory]\label{def:2d-relaxation}
Fix a unit vector $\Psi$, a declared disk exhaustion $\tpWin_m=B_{R_m}(z)$ of
the plane, or an annular exhaustion of a separately declared exterior domain,
and the operators \eqref{eq:2d-window-charge}. Infinite-volume dynamics is
formed first, fixed-window time limits second, and the spatial exhaustion
last. The three clauses are as follows.

\smallskip
\noindent\textbf{(LR1$_{\mathrm{2D}}$) Common Ces\`aro subsequence.} There is
one sequence $T_n\to\infty$ such that, for every fixed admissible window, the
two Ces\`aro states
$\omega^{\pm}_{\tpWin,n}$ formed by averaging
$\braket{\Psi,\alpha_t(A)\Psi}$ over $[T_n,2T_n]$ and $[-2T_n,-T_n]$ converge
weak-$*$, and every double-Ces\`aro average $p_{\tpWin,n}(\nu)$ of
\eqref{eq:2d-tpm} converges. By the general common-subsequence theorem of
\cref{sec:memory-index} this clause is itself a theorem --- the argument
consumes only separability of the quasi-local algebra, strong continuity of the
dynamics, and finiteness of each local spectrum, none of which sees dimension
--- but it is retained as the binder that supplies one sequence for all
clauses.

\noindent\textbf{(LR2$_{\mathrm{2D}}$) First-moment nondemolition.} With
$\tpDeph_{\tpWin,t_-}(A)=\sum_qE_{\tpWin,t_-}(\{q\})AE_{\tpWin,t_-}(\{q\})$,
the double-Ces\`aro average of
$\braket{\Psi,[\tpDeph_{\tpWin,t_-}(\Qhat_\tpWin(t_+))
-\Qhat_\tpWin(t_+)]\Psi}$ tends to zero for every fixed window along that same
sequence. This is a scalar first-moment condition, not operator asymptotic
commutativity.

\noindent\textbf{(LR3$_{\mathrm{2D}}$) First-moment tightness on the declared
exhaustion.} Writing $p_\tpWin$ for the fixed-window time limit,
\begin{equation}
  \lim_{M\to\infty}\ \sup_m\ \sum_{\abs{\nu}>M}(1+\abs{\nu})\,
  p_{\tpWin_m}(\nu)=0 .
  \label{eq:2d-lr3}
\end{equation}
One may optionally require $p_{\tpWin_m}$ to converge weakly to a probability;
that clause buys uniqueness of the limiting law and of its first moment, and is
not needed for integer support.

\smallskip
The memory observable supplied by this definition is the \emph{charge memory}
$M_Q(p):=\sum_{\nu\in\Z}\nu\,p(\nu)$, defined per subsequence, or uniquely
under the optional clause. No universal scalar displacement is defined.
\end{definition}

\provenance{M-INDEX-2D-spec, D27}{\texttt{theory/memory-index-2d.md}
steps (1.7)--(1.8)}{\texttt{theory/checks/memory\_index\_2d\_check.py}
(finite arithmetic only; it proves nothing about
(LR2$_{\mathrm{2D}}$)--(LR3$_{\mathrm{2D}}$))}%
{\texttt{theory/verdicts/w21-joint-critic.md} \S(W3)}

\begin{theorem}[Conditional two-dimensional charge ledger]%
\label{thm:memory-2d-limit}
\statusSketch\ Assume the integral-circle condition \eqref{eq:2d-int} together
with (LR1$_{\mathrm{2D}}$)--(LR3$_{\mathrm{2D}}$) of
\cref{def:2d-relaxation}. Then every ordered spatial limit point $p$ of the
window laws is a probability supported on $\Z$, and along its subsequence its
first moment obeys the charge ledger
\begin{equation}
  M_Q(p)=\lim_j\Bigl[
    \omega^{-}_{\tpWin_{m_j}}\bigl(\Qhat_{\tpWin_{m_j}}\bigr)
   -\omega^{+}_{\tpWin_{m_j}}\bigl(\Qhat_{\tpWin_{m_j}}\bigr)\Bigr].
  \label{eq:2d-ledger}
\end{equation}
Under the optional full-law convergence clause, the limit law and the value in
\eqref{eq:2d-ledger} are unique.
\end{theorem}

\begin{scope}
The conditional implication above is complete and was independently verified;
what is missing is an instance of its hypotheses. Clauses
(LR2$_{\mathrm{2D}}$) and (LR3$_{\mathrm{2D}}$) are \emph{assumed, not
derived}, and are not implied by symmetry or by a Lieb--Robinson bound; the
boundary terms they must control are perimeter-sized by
\eqref{eq:2d-perimeter}. No scalar interface displacement, channel
completeness, unconditional Goldstone instance, $1$-form-symmetry result, or
topological-order result is claimed.

The status is capped at sketch by the campaign's non-vacuity test rather than
by any defect in the argument: no nonzero model instance of
(LR2$_{\mathrm{2D}}$)--(LR3$_{\mathrm{2D}}$) has been exhibited. In
particular the exact $9\times9$ diagonalisation described below is
\emph{not} such an instance, and must not be presented as one. Promotion
requires a model in which both clauses are proved and the limiting law is not
the point mass at zero.
\end{scope}

\emph{What is proved, and what is missing.} The proved part is the implication.
At finite times, expanding the two measured values identifies the
finite-window mean as the pinched difference
\begin{equation}
  \sum_\nu\nu\,p_{\tpWin;t_-,t_+}(\nu)
   =\braket{\Qhat_\tpWin(t_-)}
    -\braket{\tpDeph_{\tpWin,t_-}\bigl(\Qhat_\tpWin(t_+)\bigr)},
  \label{eq:2d-pinched-mean}
\end{equation}
an identity in which no geometry enters. Double-Ces\`aro averaging
\eqref{eq:2d-pinched-mean} and applying (LR2$_{\mathrm{2D}}$) removes the
dephasing defect and converts the mean into the difference of the two Ces\`aro
states at each fixed window. Clause (LR3$_{\mathrm{2D}}$) supplies tightness
and uniform integrability, so Prokhorov extraction on the closed set $\Z$
yields a further weakly convergent subsequence with no loss of mass and with
the first moment passing to the limit; integer support survives because every
finite-window law already has it, the possibly drifting offset $\kappa_\tpWin$
having been cancelled at each fixed window by
\cref{thm:memory-2d-finite} before the spatial limit is taken.

The missing part is any verified instance. The relevant numerical probe is an
exact small diagonalisation of the spin-$\tfrac12$ nearest-neighbour
ferromagnet on a $9\times9$ open square, whose fully polarised state is an
exact ground state and whose one-spin-flip sector is an exact magnon sector
with the quadratic Goldstone dispersion
$\varepsilon(k)=J(2-\cos k_x-\cos k_y)$; the measured charge is the magnon
number, with on-site spectrum $\{0,1\}$. It confirms the finite theorem
convincingly --- on a $13$-site disk the law is
$p(0)=0.176188$, $p(1)=0.823812$, matching the integrated outward current to a
residue of $9\times10^{-16}$; on a $24$-site annulus, initially empty so that
magnon entry gives a negative increment, $p(-1)=0.691271$ with residue
$5\times10^{-15}$; and a seeded state straddling the disk boundary, measured at
two times with a common irrational offset $\sqrt3/11$ added to both readouts,
reproduces exact integer support with the offset cancelling as predicted. All
of that is finite-volume and finite-time. It establishes no infinite-volume
time limit and no bound of the form \eqref{eq:2d-lr3}, and the finite gap of
the $9\times9$ system is not evidence for a thermodynamic gap. Those runs are
described further in \cref{sec:numerics}.

\emph{What replaces the kink picture.} In one dimension the escaped charge
converts into a wall displacement; in two dimensions it does not, and the
honest replacement is spin-charge radiation through a closed curve. A local
disturbance launches Goldstone wave packets, and the disk record counts their
net selected charge across the circumference, an annulus recording the signed
balance across its two boundary components. Angular profiles, multipoles, and
neutral pairs are invisible to a single total-charge record. For a
zero-temperature antiferromagnet the conserved observable is a selected
\emph{uniform} spin generator; the staggered magnetisation is an order
parameter, not the conserved charge, and must not be substituted into the
theorem. Finite-window integrality remains exact in every one of these cases,
while the ordered limit remains open.

\provenance{M-INDEX-2D-spec, M-INDEX-2D-fin, LR1-GEN}%
{\texttt{theory/memory-index-2d.md} steps (1.8)--(1.11), with the conditional
theorem at step (1.9)}%
{\texttt{theory/checks/memory\_index\_2d\_check.py} (finite arithmetic only;
its green run proves nothing about (LR2$_{\mathrm{2D}}$)--(LR3$_{\mathrm{2D}}$)
and there is no gate for the limit law)}%
{\texttt{theory/verdicts/w21-joint-critic.md} \S(W3): the implication is
correct and the row is held at sketch by the non-vacuity test}
